% Standalone problem document, generated from the adjacent README.md.
% Compile with XeLaTeX. Mathematical content is maintained in Markdown.
\documentclass[11pt,a4paper]{article}
\usepackage[fontsize=10.5pt]{scrextend}
\usepackage[margin=22mm,headheight=15pt,headsep=9mm,footskip=11mm]{geometry}
\usepackage{fontspec}
\setmainfont{texgyrepagella}[Extension=.otf,UprightFont=*-regular,BoldFont=*-bold,ItalicFont=*-italic,BoldItalicFont=*-bolditalic]
\setsansfont{texgyreheros}[Extension=.otf,UprightFont=*-regular,BoldFont=*-bold,ItalicFont=*-italic,BoldItalicFont=*-bolditalic]
\setmonofont{texgyrecursor}[Extension=.otf,UprightFont=*-regular,BoldFont=*-bold,ItalicFont=*-italic,BoldItalicFont=*-bolditalic,Scale=MatchLowercase]
\usepackage{amsmath,amssymb}
\usepackage{unicode-math}
\setmathfont{texgyrepagella-math.otf}
\usepackage{xcolor,microtype,enumitem,needspace,fancyhdr}
\usepackage{bookmark}
\definecolor{ink}{HTML}{17344B}
\definecolor{accent}{HTML}{197D80}
\definecolor{muted}{HTML}{596773}
\hypersetup{colorlinks=true,linkcolor=accent,urlcolor=accent,pdftitle={MI-23 - Corrected
generalized geometric-mean product
conjecture},pdfauthor={Open Problems in Numerical Linear Algebra}}
\urlstyle{same}
\setlength{\parindent}{0pt}
\setlength{\parskip}{5pt plus 1pt minus 1pt}
\linespread{1.04}
\setlength{\emergencystretch}{3em}
\widowpenalty=10000
\clubpenalty=10000
\displaywidowpenalty=10000
\allowdisplaybreaks[1]
\setlist{nosep,leftmargin=1.5em}
\providecommand{\tightlist}{\setlength{\itemsep}{0pt}\setlength{\parskip}{0pt}}
\providecommand{\pandocbounded}[1]{#1}
\pagestyle{fancy}
\fancyhf{}
\fancyhead[L]{\sffamily\footnotesize\color{muted}OPEN PROBLEMS IN NUMERICAL LINEAR ALGEBRA}
\fancyhead[R]{\sffamily\bfseries\footnotesize\color{accent}MI-23}
\fancyfoot[L]{\parbox[b]{0.9\textwidth}{\sffamily\fontsize{8}{10}\selectfont\color{muted}Editorial ratings; a bounded search cannot certify that a problem remains unsolved.\\Literature check: 2026-09-11}}
\fancyfoot[R]{\sffamily\footnotesize\color{muted}\thepage}
\renewcommand{\headrulewidth}{0.3pt}
\renewcommand{\headrule}{\hbox to\headwidth{\color{accent}\leaders\hrule height \headrulewidth\hfill}}
\makeatletter
\renewcommand{\section}{\Needspace{5\baselineskip}\@startsection{section}{1}{0pt}{12pt plus 2pt minus 2pt}{4pt}{\sffamily\bfseries\large\color{ink}}}
\renewcommand{\subsection}{\Needspace{4\baselineskip}\@startsection{subsection}{2}{0pt}{9pt plus 1pt minus 1pt}{3pt}{\sffamily\bfseries\normalsize\color{ink}}}
\makeatother
\setcounter{secnumdepth}{0}
\begin{document}
{\sffamily\small\color{muted}Matrix inequalities and norms\par}
\vspace{3pt}
{\raggedright\hyphenpenalty=10000\exhyphenpenalty=10000\sffamily\bfseries\fontsize{21}{24}\selectfont\color{ink}Corrected
generalized geometric-mean product conjecture\par}
\vspace{7pt}
{\sffamily\small\textbf{Difficulty:} challenging\quad\textbf{Importance:} interesting
to specialist\par}
{\sffamily\small\bfseries\color{accent}Status: Solved\par}
\vspace{4pt}
{\color{accent}\hrule height 0.8pt}
\vspace{6pt}
\textbf{Rating rationale:} The four interacting exponents make this
corrected conjecture challenging; its immediate consequences concern
specialists in matrix means and log-majorization.

\subsection{Resolution --- 2026-09-11}\label{resolution-2026-09-11}

\textbf{Negative result by Matthew J. Colbrook}, Department of Applied
Mathematics and Theoretical Physics, University of Cambridge.
\textbf{Independent proof review: PASS.}

Rational positive definite \(3\times3\) matrices with \(r=s=1\), \(p=2\)
and \(t=1/8\) violate the corrected eigenvalue log-majorization
conjecture. An exact integer-power construction and rational norm
separation establish failure of the first ordered eigenvalue inequality.

The exact target is resolved. The original statement and source evidence
are retained below; its former difficulty rating is historical.

\textbf{Primary manuscript:}
\href{https://github.com/ajt60gaibb/OpenProblemsInNLA/blob/main/matrix-inequalities-and-norms/MI-23/solution.pdf}{complete
proof PDF},
\href{https://github.com/ajt60gaibb/OpenProblemsInNLA/blob/main/matrix-inequalities-and-norms/MI-23/solution.tex}{standalone
TeX}, Theorem 1.1 and its proof;
\href{https://github.com/ajt60gaibb/OpenProblemsInNLA/blob/main/matrix-inequalities-and-norms/MI-23/solution.md}{authorship
and scope}. The
\href{https://github.com/ajt60gaibb/OpenProblemsInNLA/blob/main/references/colbrook-matrix-2026-09-11/verification/reviews/MI-23-review.md}{independent
review} checks the full original argument and records its hash. The
draft was AI-assisted; this is independent agent verification, not
external human peer review or formal certification.
\href{https://github.com/ajt60gaibb/OpenProblemsInNLA/blob/main/references/colbrook-matrix-2026-09-11/README.md}{Submission
record}.

\subsection{Problem statement}\label{problem-statement}

For positive definite \(A,B\in\mathbb C^{n\times n}\), real \(r\), and
\(t\in[0,1]\), define \[
G_{r,t}(A,B)=A^{r/2}(A^{-1/2}BA^{-1/2})^tA^{r/2}.
\] Is it true, for every \(n\ge1\), every such \(A,B\), every \(p\ge1\),
\(t\in[0,1]\), and either \(r,s\ge1\) or \(r,s\le0\), that \[
\lambda\bigl(G_{r,t}(A,B)^pG_{s,1-t}(A,B)^p\bigr)
\prec_{\log}\lambda\bigl(A^{p(r+s-1)}B^p\bigr)?
\] The eigenvalues of these products are positive real numbers, ordered
decreasingly: each product is similar to a positive definite matrix. For
positive decreasing vectors \(x,y\in\mathbb R^n\), \(x\prec_{\log}y\)
means \(\prod_{j=1}^kx_j\le\prod_{j=1}^ky_j\) for \(1\le k<n\), with
equality at \(k=n\).

This compares nonlinear matrix-mean products with products of powers of
the input matrices. It would give simultaneous multiplicative control
for every partial product of ordered eigenvalues and strengthen norm
estimates used with positive definite matrix functions.

\newpage
\subsection{References}\label{references}

\begin{enumerate}
\def\labelenumi{\arabic{enumi}.}
\tightlist
\item
  M. M. Ghabries, H. Abbas, B. Mourad and A. Assi, \emph{New
  log-majorization results concerning eigenvalues and singular values
  and a complement of a norm inequality}, preprint (2021), Conjecture
  2.1, p.~6. \href{https://arxiv.org/pdf/2105.13356}{PDF}.
\item
  M. M. Ghabries, \emph{Contributions to Matrix Inequalities and Some
  Applications}, PhD thesis, University of Angers and Lebanese
  University (2022), final Open Problems, Problem 6, pp.~109--110.
  \href{https://www.researchgate.net/publication/361793582_Contributions_to_Matrix_Inequalities_and_Some_Applications}{Thesis
  uploaded by its author}.
\item
  M. M. Ghabries, \emph{A log-majorization inequality for normal
  matrices with applications to determinantal inequalities and geometric
  means} (2026), §4. \href{https://arxiv.org/html/2607.21163}{Full
  text}.
\end{enumerate}

Status check (2026-09-10): Theorem 2.1 and equation (13) of the first
source prove a substantive included region: \(r,s\ge1\), \(1\le p\le2\),
and \(\frac{r(p-1)}{(r+s)p}\le t\le\frac{rp+s}{(r+s)p}\). This includes
\(p=1\) for all \(t\). The full parameter range is the remaining
question. The first source refutes its earlier unrestricted
singular-value Conjecture 1.2; the statement here is the explicitly
proposed replacement. The thesis retains it. The July 2026 paper
concerns eigenvalue products of ordinary weighted means and does not
claim this full \(r,s,p\) statement. Current arXiv versions and searches
for the authors, ``Conjecture 2.1 generalized geometric mean'' and
2025/2026 yielded no full resolution. This is a bounded check, not
certification of openness.
\end{document}
